%=======================================================================
%              HELPER LAYER / HILFSFUNKTIONEN
%=======================================================================
% Normal users work in content.tex. This file stores content state and
% renders reusable letter components; it does not define form geometry.
% Normale Nutzer arbeiten in content.tex. Diese Datei speichert Inhalte
% und rendert wiederverwendbare Briefbausteine; sie definiert keine
% Formulargeometrie.
%
% PUBLIC: LetterBody, draft helpers, letterheading, and noticeblock.
% INTERNAL: letterhead, information, reference, and attachment rendering.
% ÖFFENTLICH: LetterBody, Entwurfshilfen, letterheading und noticeblock.
% INTERN: Briefkopf-, Informations-, Referenz- und Anlagen-Rendering.

% INTERNAL CONTENT STATE / INTERNER INHALTSZUSTAND
% Set through the public commands documented in content.tex.
% Setzen über die in content.tex dokumentierten öffentlichen Befehle.

\newcommand{\LetterSubject}{Betreff}
\newcommand{\LetterOpening}{Sehr geehrte Damen und Herren,}
\newcommand{\LetterClosing}{Mit freundlichen Grüßen}

% Optional reference fields / Optionale Referenzangaben
% Empty fields are not printed / Leere Felder werden nicht ausgegeben
% Usage in content.tex / Verwendung in content.tex:
% \SetRefAktenzeichen{...}
% \SetRefKundennummer{...}
% \SetRefIhrSchreibenVom{...}

\newcommand{\RefAktenzeichen}{}
\newcommand{\RefBGNummer}{}
\newcommand{\RefKundennummer}{}
\newcommand{\RefIhrSchreibenVom}{}
\newcommand{\RefGeschaeftszeichen}{}
\newcommand{\RefBearbeitung}{}
\newcommand{\RefIhrZeichen}{}
\newcommand{\RefIhreNachrichtVom}{}
\newcommand{\RefMeinZeichen}{}

% Optional formal number and matter data / Optionale formale Nummern- und Sachverhaltsdaten
\newcommand{\LegalGeschaeftsnummer}{}
\newcommand{\LegalSachverhalt}{}

% Default attachments by document type / Standardanlagen nach Dokumenttyp
% Used by \SetAttachmentsAuto / Verwendet von \SetAttachmentsAuto
% Custom: \SetAttachmentsCustom{...}; separate entries with \\ / Eigene Anlagen: \SetAttachmentsCustom{...}; Einträge mit \\ trennen

\newcommand{\ApplicationAttachments}{\LocalizedApplicationAttachments}
\newcommand{\AuthorityAttachments}{\LocalizedAuthorityAttachments}
\newcommand{\BusinessAttachments}{}
\newcommand{\MinimalAttachments}{}
\newcommand{\CustomAttachments}{}
\newcommand{\ResolvedAttachments}{}

%---------------------------------------------------------------
%              INTERNAL CORE HELPERS / INTERNE BASISHELFER
%---------------------------------------------------------------
% Text field check / Textfeldprüfung
% #1 text; #2 if present; #3 if empty / #1 Text; #2 wenn vorhanden; #3 wenn leer

\newcommand{\IfTextPresent}[3]{%
\ifdefempty{#1}{#3}{#2}%
}

% Page numbering / Seitenzählung
% Optional total via \UsePageTotaltrue / Optionale Gesamtseitenzahl über \UsePageTotaltrue

\newcommand{\PageNumberText}{%
\ifLanguageEnglish
\ifUsePageTotal
Page~\thepage~of~\pageref*{LastPage}%
\else
Page~\thepage%
\fi
\else
\ifUsePageTotal
Seite~\thepage~von~\pageref*{LastPage}%
\else
Seite~\thepage%
\fi
\fi
}

% Contact separator / Kontakttrenner
% Between contact details in the letterhead / Zwischen Kontaktangaben im Briefkopf

\newcommand{\ContactSeparator}{%
\quad\textcolor{linegray}{|}\quad
}

% Optional icon output / Optionale Icon-Ausgabe
% No icons in authority mode / Keine Icons im Behördenmodus
% #1 FontAwesome icon; #2 text / #1 FontAwesome-Icon; #2 Text

\newcommand{\MaybeIcon}[2]{%
\ifdefstring{\DocumentMode}{authority}{%
#2%
}{%
\ifUseIcons
\ifFontAwesomeAvailable
\faIcon{#1}~#2%
\else
#2%
\fi
\else
#2%
\fi
}%
}

%---------------------------------------------------------------
%                Letterhead / Briefkopf
%---------------------------------------------------------------
% Contact details / Kontaktdaten
% Values from data/sender-data.tex / Werte aus data/sender-data.tex

\newcommand{\PhoneContact}{\MaybeIcon{phone}{\LabelPhone~\SenderPhone}}
\newcommand{\EmailContact}{\MaybeIcon{envelope}{\href{mailto:\SenderEmail}{\SenderEmail}}}
\newcommand{\LinkedInContact}{\MaybeIcon{linkedin}{\href{\SenderLinkedInUrl}{LinkedIn}}}
\newcommand{\WebsiteContact}{\MaybeIcon{globe}{\href{\SenderWebsiteUrl}{\SenderWebsiteText}}}

% Second contact line / Zweite Kontaktzeile
% Optional LinkedIn and website / Optional LinkedIn und Website

\newcommand{\MaybeOnlineContactRow}{%
\ifUseLinkedIn
\\[0.7mm]{\fontsize{\HeaderContactSize}{\HeaderContactLeading}\selectfont
\LinkedInContact
\ifUseWebsite\ContactSeparator\WebsiteContact\fi}%
\else
\ifUseWebsite
\\[0.7mm]{\fontsize{\HeaderContactSize}{\HeaderContactLeading}\selectfont
\WebsiteContact}%
\fi
\fi
}

% Position or profile line / Positions- oder Profilzeile
% Requires \UsePositiontrue and \SenderPosition / Benötigt \UsePositiontrue und \SenderPosition

\newcommand{\MaybePositionRow}{%
\ifUsePosition
\IfTextPresent{\SenderPosition}{%
{\fontsize{\HeaderPositionSize}{\HeaderPositionLeading}\selectfont
\textbf{\textcolor{accent}{\SenderPosition}}}\\[1.6mm]%
}{}%
\fi
}

% Left letterhead content / Linker Briefkopf
% Name, position, and contact details / Name, Position und Kontaktdaten

\newcommand{\HeaderLeftContent}{%
{\fontsize{\HeaderNameSize}{\HeaderNameLeading}\selectfont
\textcolor{textlight}{\SenderFirstName}~\textbf{\textcolor{textdark}{\SenderLastName}}
}\\[-0.1mm]
\MaybePositionRow
{\fontsize{\HeaderContactSize}{\HeaderContactLeading}\selectfont
\PhoneContact
\ContactSeparator
\EmailContact
}%
\MaybeOnlineContactRow
}

% Image logo / Bildlogo
% Uses \HeaderLogoFile; monogram fallback if the selected file is missing.
% Nutzt \HeaderLogoFile; Monogramm-Fallback, falls die gewählte Datei fehlt.

\newcommand{\ImageLogo}{%
\IfFileExists{\HeaderLogoFile}{%
  \makebox[0pt][r]{%
    \raisebox{-\height}[0pt][0pt]{%
      \includegraphics[width=50mm,height=\LetterFormLogoHeight,keepaspectratio]{\HeaderLogoFile}%
    }%
  }%
}{\MonogramLogo}%
}

% Monogram fallback / Monogramm-Fallback
% Minimal mode uses a reduced black “mm” mark / Minimalmodus nutzt ein reduziertes schwarzes „mm“-Zeichen

\newcommand{\MonogramLogo}{%
\def\MonogramColor{accent}%
\def\MonogramText{MM}%
\ifdefstring{\DocumentMode}{minimal}{%
  \def\MonogramColor{black}%
  \def\MonogramText{MM}%
}{}%
\begin{tikzpicture}
\node[
circle,
draw=\MonogramColor,
line width=0.6pt,
fill=none,
minimum width=10mm,
minimum height=10mm,
inner sep=0pt
] {\fontsize{\LetterFormMonogramSize}{\LetterFormMonogramSize}\selectfont\bfseries\textcolor{\MonogramColor}{\MonogramText}};
\end{tikzpicture}%
}

% Draft warning in the main letter flow. \DraftMode is authoritative;
% \UseDraftNotice... is retained only for source compatibility.
\newcommand{\DraftNoticeInLetterFlow}{%
  \ifDraftMode
    \par\noindent\textcolor{red}{\textbf{\LabelDraftNotice}}\par\medskip
  \fi
}

%---------------------------------------------------------------
%     DEMO-DATA REMINDER / HINWEIS AUF DEMO-DATEN
%---------------------------------------------------------------
% When draft mode is OFF and the real entry point (content.tex) is compiled,
% the log warns if shipped demo values are still present. Log-only; this is
% never printed onto the letter. Disable with \CheckDemoDatafalse.
% Wenn der Entwurfsmodus AUS ist und content.tex kompiliert wird, warnt der
% Log-Eintrag bei verbliebenen Demo-Werten. Nur im Log, nie im Brief sichtbar.
% Mit \CheckDemoDatafalse abschalten.

% Captured demo recipient default from content.tex; the group keeps the exact
% argument tokenization for an \ifx comparison without touching \RecipientBlock.
% Erfasster Demo-Empfänger aus content.tex; die Gruppe bewahrt die exakte
% Tokenfolge für den \ifx-Vergleich, ohne \RecipientBlock zu verändern.
\newcommand{\DemoRecipientDefault}{}
\begingroup
\SetManualRecipient{%
Beispielorganisation\\
Allgemeine Korrespondenz\\
Musterweg 1\\
12345 Beispielstadt
}%
\global\let\DemoRecipientDefault\RecipientBlock
\endgroup

\newcommand{\WarnDemoField}[1]{%
  \PackageWarning{letter-template}{%
    #1 still contains demo data; replace it before sending, or set
    \string\CheckDemoDatafalse\space to silence this check%
  }%
}

\newcommand{\RunDemoDataCheck}{%
  \ifdefstring{\SenderName}{Max Mustermann}%
    {\WarnDemoField{Sender name (data/sender-data.tex)}}{}%
  \ifdefstring{\LetterSubject}{Betreff des eigenen Schreibens}%
    {\WarnDemoField{Subject (content.tex)}}{}%
  \ifdefstring{\RecipientMode}{manual}{%
    \ifdefequal{\RecipientBlock}{\DemoRecipientDefault}%
      {\WarnDemoField{Recipient (content.tex)}}{}%
  }{}%
}

% Gated to the real entry point only: content.tex compiled as jobname "main".
% Guide previews via \ContentFile (jobname main) and script guide builds
% (jobname = guide name) both stay silent.
% Nur am echten Einstieg aktiv: content.tex mit Jobname "main". Guide-Vorschauen
% über \ContentFile und Skript-Guidebuilds bleiben still.
\AtEndDocument{%
  \ifCheckDemoData
    \ifDraftMode\else
      \edef\DemoCheckJob{\jobname}%
      \ifdefstring{\ContentFile}{content.tex}{%
        \ifdefstring{\DemoCheckJob}{main}{\RunDemoDataCheck}{}%
      }{}%
    \fi
  \fi
}

% TikZ symbol; no external image required / TikZ-Symbol; keine externe Bilddatei erforderlich

\newcommand{\TikzLogo}{%
\begin{tikzpicture}[x=1mm,y=1mm]
\draw[line width=0.9pt, color=accent, line cap=round]
(0,4) .. controls (7,13) and (18,13) .. (25,4);
\draw[line width=0.7pt, color=textdark, line cap=round]
(4,2.5) -- (21,2.5);
\draw[line width=0.7pt, color=accent, line cap=round]
(12.5,2.5) -- (12.5,10.5);
\fill[accent] (12.5,10.5) circle (1);
\end{tikzpicture}%
}

% Header variant selection / Auswahl der Briefkopfvariante
% Possible values / Mögliche Werte:
% image = logo.png; colorblock = decorative color block / dekorativer Farbblock;
% monogram = initials / Initialen; tikz = TikZ symbol
% greenblock remains a compatibility alias for colorblock / greenblock bleibt ein Kompatibilitätsalias für colorblock

\newcommand{\IfColorBlockHeaderTF}[2]{%
\ifdefstring{\HeaderVariant}{colorblock}{#1}{%
\ifdefstring{\HeaderVariant}{greenblock}{#1}{#2}%
}%
}

\newcommand{\BrandLogo}{%
\ifdefstring{\HeaderVariant}{image}{\ImageLogo}{%
\ifdefstring{\HeaderVariant}{tikz}{\TikzLogo}{\MonogramLogo}%
}%
}

% Right header logo / Logo rechts im Briefkopf
% Selected logo, right-aligned / Gewähltes Logo, rechtsbündig

\newcommand{\HeaderLogoContent}{\raggedleft\ifUseHeaderLogo\BrandLogo\fi}

% Business letterhead: logo at the active form profile position.
% Geschäftsbriefkopf: Logo an der Position des aktiven Formprofils.
\newcommand{\BusinessHeaderLogo}{%
\IfFileExists{\HeaderLogoFile}{%
  \includegraphics[width=75mm,height=\LetterFormLogoHeight,keepaspectratio]{\HeaderLogoFile}%
}{\MonogramLogo}%
}

\newcommand{\BusinessHeaderContent}{%
\ifUseHeaderLogo
\begin{tikzpicture}[remember picture,overlay]
  \node[anchor=north west,inner sep=0pt,outer sep=0pt]
    at ([xshift=\LetterFormBusinessLogoX,yshift=-\LetterFormBusinessLogoTop]current page.north west)
    {\makebox[75mm][c]{\BusinessHeaderLogo}};
\end{tikzpicture}%
\fi
}

%---------------------------------------------------------------
%          Business information block / Informationsblock
%---------------------------------------------------------------
% Shared 75 mm component; its page position comes from the active form profile.
% Gemeinsamer 75-mm-Baustein; seine Seitenposition kommt aus dem aktiven Formprofil.
% Empty optional values are omitted / Leere optionale Werte entfallen

\newcommand{\BusinessInformationRow}[2]{%
  {\raggedright\textcolor{textmid}{#1}} & {\raggedright #2}\tabularnewline[0.25mm]
}

\newcommand{\MaybeBusinessInformationRow}[3]{%
\IfTextPresent{#2}{\BusinessInformationRow{#1}{#3}}{}%
}

% Allow a deliberate line break after @ in the shared 75 mm information
% block while keeping the regular text font.
\def\BusinessEmailWithBreak#1@#2\BusinessEmailEnd{#1@\allowbreak#2}
\newcommand{\BusinessEmailText}{%
  \expandafter\BusinessEmailWithBreak\SenderEmail\BusinessEmailEnd
}

\newcommand{\BusinessInformationBlock}{%
\begingroup
\fontsize{11pt}{14.3pt}\selectfont
\setlength{\tabcolsep}{0pt}%
\begin{tabular}{@{}p{20mm}@{\hspace{1.5mm}}p{53.5mm}@{}}
\MaybeBusinessInformationRow{\LabelBusinessContactPerson}{\SenderContactPerson}{\SenderContactPerson}%
\MaybeBusinessInformationRow{\LabelBusinessDepartment}{\SenderDepartment}{\SenderDepartment}%
\MaybeBusinessInformationRow{\LabelBusinessPhone}{\SenderPhone}{\SenderPhone}%
\MaybeBusinessInformationRow{\LabelBusinessFax}{\SenderFax}{\SenderFax}%
\MaybeBusinessInformationRow{\LabelBusinessEmail}{\SenderEmail}{\href{mailto:\SenderEmail}{\BusinessEmailText}}%
\MaybeBusinessInformationRow{\LabelBusinessWebsite}{\SenderWebsiteText}{\href{\SenderWebsiteUrl}{\SenderWebsiteText}}%
\BusinessInformationRow{\LabelBusinessDate}{\today}%
\end{tabular}
\endgroup
}

% The stored value is evaluated when KOMA renders the letter / Der gespeicherte Wert wird beim Setzen des Briefs ausgewertet
\newcommand{\BusinessLocationContent}{%
\ifdefstring{\DocumentMode}{business}{\BusinessInformationBlock}{}%
}

% Authority-only reference block at the active profile position.
% Behörden-Referenzblock an der Position des aktiven Formprofils.
% The shared component totals 75 mm; empty optional rows are omitted.
% Der gemeinsame Baustein ist 75 mm breit; leere optionale Zeilen entfallen.

\newcommand{\AuthorityReferenceRow}[2]{%
  {\raggedright\bfseries\textcolor{textmid}{#1}\par} & {\raggedright #2}\tabularnewline[0.25mm]
}

\newcommand{\MaybeAuthorityReferenceRow}[2]{%
\IfTextPresent{#2}{\AuthorityReferenceRow{#1}{#2}}{}%
}

% Prefer letter-specific \Ref... data and otherwise use neutral authority defaults / Briefspezifische \Ref...-Daten bevorzugen, sonst neutrale Behörden-Standardwerte nutzen
\newcommand{\AuthorityFallbackRow}[3]{%
\ifdefempty{#2}{%
  \ifdefempty{#3}{}{\AuthorityReferenceRow{#1}{#3}}%
}{\AuthorityReferenceRow{#1}{#2}}%
}

\newcommand{\AuthorityTwoLineFallbackRow}[4]{%
\ifdefempty{#2}{%
  \ifdefempty{#3}{%
    \ifdefempty{#4}{}{\AuthorityReferenceRow{#1}{#4}}%
  }{%
    \AuthorityReferenceRow{#1}{#3\ifdefempty{#4}{}{\\#4}}%
  }%
}{\AuthorityReferenceRow{#1}{#2}}%
}

\newcommand{\AuthorityReferenceBlock}{%
\begingroup
\fontsize{11pt}{14.3pt}\selectfont
\setlength{\tabcolsep}{0pt}%
\begin{tabular}{@{}p{39mm}@{\hspace{1.5mm}}p{34.5mm}@{}}
\AuthorityTwoLineFallbackRow{\LabelRefIhrZeichen}{\RefIhrZeichen}{\AuthorityIhrZeichenLineOne}{\AuthorityIhrZeichenLineTwo}%
\AuthorityFallbackRow{\LabelRefIhreNachrichtVom}{\RefIhreNachrichtVom}{\AuthorityIhreNachrichtVom}%
\AuthorityTwoLineFallbackRow{\LabelRefMeinZeichen}{\RefMeinZeichen}{\AuthorityMeinZeichenLineOne}{\AuthorityMeinZeichenLineTwo}%
\AuthorityFallbackRow{\LabelRefKundennummer}{\RefKundennummer}{\AuthorityKundennummer}%
\AuthorityFallbackRow{\LabelRefBGNummer}{\RefBGNummer}{\AuthorityBGNummer}%
\MaybeAuthorityReferenceRow{\LabelBusinessFax}{\SenderFax}%
\AuthorityReferenceRow{\LabelRefDatum}{\AuthorityDatum}%
\end{tabular}
\endgroup
}

\newcommand{\AuthorityLocationContent}{%
\ifdefstring{\DocumentMode}{authority}{\AuthorityReferenceBlock}{}%
}

% Optional formal matter block in the area reserved by the active profile.
% Optionaler formaler Sachverhaltsblock im Bereich des aktiven Formprofils.
% Compact outlined fields; empty fields are omitted.
% Kompakte Rahmenfelder; leere Felder werden nicht ausgegeben.
\newcommand{\LegalCaseField}[2]{%
  \begingroup
  \setlength{\fboxsep}{2mm}%
  \setlength{\fboxrule}{0.6pt}%
  \fcolorbox{linegray}{white}{%
    \parbox[t]{65mm}{%
      \raggedright
      {\fontsize{7pt}{8.5pt}\selectfont\textcolor{textmid}{#1}\par}%
      \vspace{1.1mm}%
      {\fontsize{9pt}{11pt}\selectfont\bfseries #2\strut}%
    }%
  }%
  \endgroup
}

\newcommand{\MaybeLegalCaseField}[2]{%
  \IfTextPresent{#2}{\LegalCaseField{#1}{#2}\par\vspace{1.8mm}}{}%
}

\newcommand{\LegalCaseBlock}{%
\normalfont
\MaybeLegalCaseField{\LabelLegalGeschaeftsnummer}{\LegalGeschaeftsnummer}%
\MaybeLegalCaseField{\LabelLegalSachverhalt}{\LegalSachverhalt}%
}

\newcommand{\LegalLocationContent}{%
\ifdefstring{\DocumentMode}{legal}{%
  \ifdefempty{\LegalGeschaeftsnummer}{%
    \ifdefempty{\LegalSachverhalt}{}{\LegalCaseBlock}%
  }{\LegalCaseBlock}%
}{}%
}

% Formal horizontal reference block above the subject.
% Horizontaler formaler Referenzblock oberhalb des Betreffs.
\newcommand{\LegalReferenceColumn}[2]{%
\parbox[t]{38mm}{%
  \raggedright
  \normalfont\fontsize{8pt}{10pt}\selectfont
  \textbf{\textcolor{textmid}{#1}}\par
  \vspace{1mm}%
  #2%
}%
}

\newcommand{\LegalFallbackColumn}[3]{%
\ifdefempty{#2}{%
  \ifdefempty{#3}{}{\LegalReferenceColumn{#1}{#3}}%
}{\LegalReferenceColumn{#1}{#2}}%
}

\newcommand{\LegalTwoLineFallbackColumn}[4]{%
\ifdefempty{#2}{%
  \ifdefempty{#3}{%
    \ifdefempty{#4}{}{\LegalReferenceColumn{#1}{#4}}%
  }{\LegalReferenceColumn{#1}{#3\ifdefempty{#4}{}{\\#4}}}%
}{\LegalReferenceColumn{#1}{#2}}%
}

\newcommand{\LegalReferenceBlock}{%
\begingroup
\normalfont
\makebox[\linewidth][l]{%
  \LegalTwoLineFallbackColumn{\LabelRefIhrZeichen}{\RefIhrZeichen}{\AuthorityIhrZeichenLineOne}{\AuthorityIhrZeichenLineTwo}%
  \hfill
  \LegalTwoLineFallbackColumn{\LabelRefMeinZeichen}{\RefMeinZeichen}{\AuthorityMeinZeichenLineOne}{\AuthorityMeinZeichenLineTwo}%
  \hfill
  \LegalFallbackColumn{\LabelRefIhreNachrichtVom}{\RefIhreNachrichtVom}{\AuthorityIhreNachrichtVom}%
  \hfill
  \LegalReferenceColumn{\LabelRefOrtDatum}{\SenderPlace, \today}%
}%
\endgroup
}

\newcommand{\LegalSubjectContent}{%
\LegalReferenceBlock
\par\vspace{3.6mm}%
\LetterSubject
}

\newcommand{\ResolvedSubjectContent}{%
\ifdefstring{\DocumentMode}{legal}{\LegalSubjectContent}{\LetterSubject}%
}

%---------------------------------------------------------------
%                    Reference block / Referenzblock
%---------------------------------------------------------------
% Content flag / Inhaltsmerker

\newif\ifHasReferenceContent

% Reference field check / Referenzfeldprüfung
% Enables the block if any field is set / Aktiviert den Block, sobald ein Feld gesetzt ist

\newcommand{\CheckReferenceContent}{%
\HasReferenceContentfalse
\ifdefempty{\RefAktenzeichen}{}{\HasReferenceContenttrue}%
\ifdefempty{\RefBGNummer}{}{\HasReferenceContenttrue}%
\ifdefempty{\RefKundennummer}{}{\HasReferenceContenttrue}%
\ifdefempty{\RefIhrSchreibenVom}{}{\HasReferenceContenttrue}%
\ifdefempty{\RefGeschaeftszeichen}{}{\HasReferenceContenttrue}%
\ifdefempty{\RefBearbeitung}{}{\HasReferenceContenttrue}%
}

% Reference row / Referenzzeile
% #1 label; #2 value / #1 Bezeichnung; #2 Wert

\newcommand{\ReferenceRow}[2]{%
\textcolor{textmid}{#1} & #2\\[-0.2mm]
}

% Reference block content / Referenzblockinhalt
% Empty fields are skipped / Leere Felder werden übersprungen

\newcommand{\ReferenceBlockContent}{%
\begingroup
\fontsize{8.6pt}{10.8pt}\selectfont
\setlength{\tabcolsep}{0pt}
\begin{tabular}{@{}p{31mm}p{94mm}@{}}
\ifdefempty{\RefAktenzeichen}{}{\ReferenceRow{\LabelRefAktenzeichen}{\RefAktenzeichen}}%
\ifdefempty{\RefBGNummer}{}{\ReferenceRow{\LabelRefBGNummer}{\RefBGNummer}}%
\ifdefempty{\RefKundennummer}{}{\ReferenceRow{\LabelRefKundennummer}{\RefKundennummer}}%
\ifdefempty{\RefIhrSchreibenVom}{}{\ReferenceRow{\LabelRefIhrSchreibenVom}{\RefIhrSchreibenVom}}%
\ifdefempty{\RefGeschaeftszeichen}{}{\ReferenceRow{\LabelRefGeschaeftszeichen}{\RefGeschaeftszeichen}}%
\ifdefempty{\RefBearbeitung}{}{\ReferenceRow{\LabelRefBearbeitung}{\RefBearbeitung}}%
\end{tabular}
\endgroup
}

% Reference block output / Referenzblockausgabe
% Requires enabled mode and content / Benötigt aktivierten Modus und Inhalt

\newcommand{\MaybeReferenceBlock}{%
\ifdefstring{\DocumentMode}{authority}{}{%
\ifUseReferenceBlock
\CheckReferenceContent
\ifHasReferenceContent
\par\vspace{0.2em}
\noindent{\color{linegray}\rule{\linewidth}{0.35pt}}%
\par\vspace{0.35em}
\ReferenceBlockContent
\par\vspace{0.55em}
{\color{linegray}\rule{\linewidth}{0.35pt}}%
\par\vspace{1.0em}
\fi
\fi
}%
}

%---------------------------------------------------------------
%                    Attachments / Anlagen
%---------------------------------------------------------------
% Possible modes / Mögliche Modi:
% auto = mode defaults / Modusstandard; custom = custom / eigene; none = disabled / deaktiviert

\newcommand{\ResolveAttachments}{%
\renewcommand{\ResolvedAttachments}{}%
\ifdefstring{\AttachmentMode}{none}{%
\UseAttachmentsfalse
}{%
\ifdefstring{\AttachmentMode}{custom}{%
\let\ResolvedAttachments\CustomAttachments
}{%
\ifdefstring{\AttachmentPreset}{application}{%
\let\ResolvedAttachments\ApplicationAttachments
}{%
\ifdefstring{\AttachmentPreset}{authority}{%
\let\ResolvedAttachments\AuthorityAttachments
}{%
\ifdefstring{\AttachmentPreset}{business}{%
\let\ResolvedAttachments\BusinessAttachments
}{%
\let\ResolvedAttachments\MinimalAttachments
}%
}%
}%
}%
\ifdefempty{\ResolvedAttachments}{\UseAttachmentsfalse}{\UseAttachmentstrue}%
}%
}

% Attachments output / Anlagenausgabe
% Printed only if enabled and not empty / Nur aktiv und nicht leer ausgegeben

\newcommand{\MaybeAttachments}{%
\ResolveAttachments
\ifUseAttachments
\ifdefempty{\ResolvedAttachments}{}{\encl{\ResolvedAttachments}}%
\fi
}

%---------------------------------------------------------------
%             PUBLIC: DRAFT HELPERS / ÖFFENTLICH: ENTWURFSHILFEN
%---------------------------------------------------------------
% Visible only with \DraftModetrue / Nur mit \DraftModetrue sichtbar

\newcommand{\DraftTag}[2]{%
\ifDraftMode
\begingroup
\setlength{\fboxsep}{2.5pt}%
\colorbox{draftyellow}{%
{\fontsize{8pt}{9.5pt}\selectfont
\textcolor{draftorange}{\textbf{#1:} #2}}%
}%
\endgroup
\fi
}

% TODO marker / TODO-Markierung
% Usage / Nutzung:
% \TODO{Check text / Text prüfen}

\newcommand{\TODO}[1]{\DraftTag{TODO}{#1}}

% Draft note / Entwurfsnotiz
% Usage / Nutzung:
% \DRAFTNOTE{Add follow-up question / Rückfrage ergänzen}

\newcommand{\DRAFTNOTE}[1]{\DraftTag{\LabelDraftNote}{#1}}

% Placeholder / Platzhalter
% Usage / Nutzung:
% \PLACEHOLDER{Add paragraph / Absatz ergänzen}

\newcommand{\PLACEHOLDER}[1]{\ifDraftMode\textcolor{draftorange}{[#1]}\fi}

%---------------------------------------------------------------
%              PUBLIC: TEXT BLOCKS / ÖFFENTLICH: TEXTBAUSTEINE
%---------------------------------------------------------------
% Usage / Nutzung:
% \letterheading{Heading / Überschrift}

\newcommand{\letterheading}[1]{%
\vspace{0.75em}
{\normalsize\bfseries\textcolor{textdark}{#1}}%
\hspace{0.55em}{\color{linegray}\rule[0.45ex]{0.34\linewidth}{0.4pt}}%
\par\vspace{0.25em}
}

% Notice block / Hinweisblock
% Usage / Nutzung:
% \noticeblock{Notice text / Hinweistext}

\newcommand{\noticeblock}[1]{%
\vspace{0.6em}
\begingroup
\setlength{\fboxsep}{7pt}
\noindent
\colorbox{noticebackground}{%
\parbox{\dimexpr\linewidth-2\fboxsep\relax}{%
\fontsize{9.5pt}{12pt}\selectfont #1%
}%
}
\endgroup
\vspace{0.4em}
}

%---------------------------------------------------------------
%       PUBLIC: NATURAL LETTER BODY / ÖFFENTLICH: BRIEFTEXT
%---------------------------------------------------------------
% Wraps the technical scrlttr2 letter structure so users can write normal
% paragraphs between \begin{LetterBody} and \end{LetterBody}.
% Kapselt die technische Briefstruktur, damit Nutzer normalen Fließtext
% zwischen \begin{LetterBody} und \end{LetterBody} schreiben können.

\newenvironment{LetterBody}{%
  \begin{letter}{\ResolvedRecipientBlock}%
  \opening{\MaybeReferenceBlock\DraftNoticeInLetterFlow\LetterOpening}%
}{%
  \closing{\LetterClosing}%
  \MaybeAttachments
  \end{letter}%
}

